\documentclass[11pt,aspectratio=169]{beamer}

\usepackage[utf8]{inputenc}
\usepackage[T1]{fontenc}
\usepackage[french]{babel}
\shorthandoff{:}
\usepackage{csquotes}
\usepackage{libertinus}
\usepackage{microtype}
\usepackage{amsmath,amssymb,mathtools}
\usepackage{booktabs,tabularx,array,colortbl}
\usepackage{graphicx}
\makeatletter
% Les sources longues portent des descriptions alternatives PDF. Beamer ne
% prend pas encore en charge DocumentMetadata : la clé reste donc sans effet
% dans le diaporama, sans qu'il faille dupliquer les appels d'image.
\define@key{Gin}{alt}{}
\makeatother
\usepackage[export]{adjustbox}
\usepackage{tikz}
\usetikzlibrary{arrows.meta,positioning,calc,decorations.pathreplacing,fit,backgrounds,patterns}
\tikzset{alt/.code={}}
\usepackage{pgfplots}
\usepgfplotslibrary{statistics}
\pgfplotsset{compat=1.18}
\usepackage{xcolor}
\usepackage{xspace}
\usepackage{listings}
\usepackage[backend=biber,style=numeric,sorting=none]{biblatex}
\addbibresource{references.bib}
\AtEveryBibitem{\clearfield{urldate}}

\definecolor{ink}{HTML}{18212B}
\definecolor{accent}{HTML}{155E75}
\definecolor{accentlight}{HTML}{E7F3F6}
\definecolor{warning}{HTML}{9A3412}
\definecolor{warninglight}{HTML}{FFF1E8}
\definecolor{gridgray}{HTML}{CBD5E1}
\definecolor{rulecolor}{HTML}{4D7C0F}
\definecolor{rulelight}{HTML}{F1F8E8}
\definecolor{gryff}{HTML}{155E75}
\definecolor{huffle}{HTML}{9A3412}
\definecolor{raven}{HTML}{4D7C0F}
\definecolor{slyth}{HTML}{6B21A8}
\definecolor{herbo}{HTML}{9D174D}
\definecolor{runes}{HTML}{1D4ED8}
\definecolor{courseproblem}{HTML}{687386}
\definecolor{coursedata}{HTML}{47777B}
\definecolor{coursescore}{HTML}{356B8B}
\definecolor{courseprobability}{HTML}{5B6093}
\definecolor{courseloss}{HTML}{8A5D68}
\definecolor{courseoptimization}{HTML}{98633F}
\definecolor{courseevaluation}{HTML}{527460}

% Grille 16:9 : toute la surface sert au cours, sans chrome de navigation.
\renewcommand{\familydefault}{\sfdefault}
\usefonttheme{professionalfonts}
\newlength{\coursesidemargin}
\setlength{\coursesidemargin}{7mm}
\setbeamersize{text margin left=\coursesidemargin,text margin right=\coursesidemargin}
\setbeamertemplate{navigation symbols}{}
\setbeamertemplate{headline}{}
\setbeamertemplate{footline}{}
\setbeamertemplate{caption}[numbered]
\setbeamertemplate{frametitle continuation}{%
  \ifnum\insertcontinuationcount>1
    \space{\small\mdseries\color{white!62}(suite)}%
  \fi}
\setbeamercolor{normal text}{fg=ink,bg=white}
\setbeamercolor{structure}{fg=accent}
\setbeamercolor{frametitle}{fg=white,bg=ink}
\setbeamercolor{caption name}{fg=accent!72!ink}
\setbeamercolor{caption}{fg=ink!62}
\setbeamerfont{normal text}{size=\normalsize}
\setbeamerfont{frametitle}{size=\footnotesize,series=\bfseries}
\setbeamerfont{framecrumb}{size=\tiny,series=\bfseries}
\setbeamerfont{caption}{size=\scriptsize}
\setbeamerfont{section title}{size=\LARGE,series=\bfseries}
\setbeamerfont{part title}{size=\LARGE,series=\bfseries}
\setbeamertemplate{frametitle}{%
  \nointerlineskip
  \begin{beamercolorbox}[wd=\paperwidth,ht=3.8mm,dp=1.2mm]{frametitle}%
    \hbox to \paperwidth{%
      \parbox[c][5mm][c]{.20\paperwidth}{%
        \hspace*{4mm}{\usebeamerfont{framecrumb}\color{white!68}%
          \MakeUppercase{\coursepartposition}}}%
      \parbox[c][5mm][c]{.60\paperwidth}{%
        \centering
        \begin{adjustbox}{max width=.58\paperwidth}%
          {\usebeamerfont{frametitle}\color{white}%
            \MakeUppercase{\insertframetitle}}%
        \end{adjustbox}}%
      \parbox[c][5mm][c]{.20\paperwidth}{}%
    }%
  \end{beamercolorbox}%
  \nointerlineskip
  \coursenavigation}
\setbeamertemplate{itemize item}{\color{accent!62}\small\textbullet}
\setbeamertemplate{itemize subitem}{\color{accent!48}\scriptsize\textbullet}
\setbeamertemplate{enumerate item}{%
  \colorbox{accent!12}{\color{accent!82!ink}\scriptsize\bfseries\insertenumlabel}}
\setlength{\parskip}{.42em}
\setlength{\emergencystretch}{2em}
\linespread{1.05}

\title{Régression logistique}
\subtitle{Modèle, optimisation et évaluation}
\author{Sylvain Maitre}
\date{Version 2.0}
\setbeamertemplate{title page}{%
  \vfill
  {\fontsize{29}{33}\selectfont\bfseries\color{ink}\inserttitle\par}
  \vspace{3mm}
  {\fontsize{16}{20}\selectfont\color{accent}\insertsubtitle\par}
  \vspace{7mm}{\color{accent!42}\rule{31mm}{1pt}}\par
  \vfill
  {\small\color{ink!72}\insertauthor\par}
  \vspace{1mm}
  {\scriptsize\color{ink!48}\insertdate\par}
  \vspace{4mm}}

% Les images respectent simultanément la largeur disponible et la hauteur utile.
\let\courseincludegraphics\includegraphics
\RenewDocumentCommand{\includegraphics}{O{}m}{%
  \courseincludegraphics[max height=.67\textheight,keepaspectratio,#1]{#2}}
\NewDocumentEnvironment{slidevisual}{}{%
  \begin{adjustbox}{max totalsize={\textwidth}{.75\textheight},center}%
  \begin{minipage}{\textwidth}}
  {\end{minipage}\end{adjustbox}}

\pgfplotsset{courseplot/.style={width=.76\textwidth,height=5.25cm,
  grid=major,grid style={gridgray!40}}}
\pgfplotsset{cloud/.style={width=.72\textwidth,height=5.4cm,
  axis lines=box,axis line style={ink!35},
  grid=major,grid style={gridgray!35},
  tick label style={font=\scriptsize},label style={font=\scriptsize},
  xlabel style={yshift=-2pt},ylabel style={yshift=4pt},
  legend style={at={(0.5,1.08)},anchor=south,legend columns=-1,draw=none,
    font=\scriptsize,/tikz/every even column/.append style={column sep=8pt}}}}
\pgfplotsset{panel/.style={width=.3\textwidth,height=4.15cm,
  axis lines=box,axis line style={ink!35},
  grid=major,grid style={gridgray!35},
  tick label style={font=\tiny},label style={font=\tiny},
  title style={font=\scriptsize,yshift=-2pt},
  xmin=-2.8,xmax=2.8,ymin=-2.8,ymax=2.8,
  xtick={-2,0,2},ytick={-2,0,2}}}

\lstset{
  basicstyle=\ttfamily\scriptsize,
  keywordstyle=\color{accent}\bfseries,
  commentstyle=\color{rulecolor},
  stringstyle=\color{warning},
  showstringspaces=false,
  breaklines=true,
  breakatwhitespace=true,
  frame=none,
  numbers=none,
  columns=fullflexible,
  language=Python,
  literate={é}{{\'e}}1 {è}{{\`e}}1 {ê}{{\^e}}1
    {à}{{\`a}}1 {ç}{{\c c}}1}
\lstdefinestyle{code}{
  basicstyle=\ttfamily\scriptsize,
  xleftmargin=3mm,
  aboveskip=.5em,
  belowskip=.4em}

% Numérotation stable entre le cours long et le diaporama.
\newcounter{coursechapter}
\newcommand{\coursechapterkind}{Chapitre}
\newcommand{\coursechapterindicator}{}
\newcommand{\coursepartshort}{Vue d'ensemble}
\newcommand{\coursepartposition}{Vue d'ensemble}
\newcommand{\coursepartindex}{0}
\newcommand{\coursepartcolor}{courseproblem}
\newcommand{\coursestage}[3]{%
  \begingroup
  \ifnum#1=\coursepartindex
    \def\stageforeground{#3}%
    \def\stagemarker{\textbullet\space}%
    \def\stageweight{\bfseries}%
  \else
    \ifnum#1<\coursepartindex
      \def\stageforeground{ink!52}%
    \else
      \def\stageforeground{ink!18}%
    \fi
    \def\stagemarker{}%
    \def\stageweight{\mdseries}%
  \fi
  \parbox[c][3mm][c]{.142857\paperwidth}{%
    \centering
    {\fontsize{5.55}{6}\selectfont\stageweight
      \color{\stageforeground}\stagemarker\MakeUppercase{#2}}}%
  \endgroup}
\newcommand{\coursenavigation}{%
  \hbox to \paperwidth{%
    \coursestage{1}{Problème}{courseproblem}%
    \coursestage{2}{Données}{coursedata}%
    \coursestage{3}{Score}{coursescore}%
    \coursestage{4}{Probabilité}{courseprobability}%
    \coursestage{5}{Perte}{courseloss}%
    \coursestage{6}{Optimisation}{courseoptimization}%
    \coursestage{7}{Évaluation}{courseevaluation}}}
\renewcommand{\thecoursechapter}{\arabic{coursechapter}}
\newcommand{\coursereperes}{%
  \setcounter{coursechapter}{0}%
  \renewcommand{\coursechapterkind}{Repère}%
  \renewcommand{\thecoursechapter}{\Roman{coursechapter}}}
\newcommand{\coursemain}{%
  \setcounter{coursechapter}{0}%
  \renewcommand{\coursechapterkind}{Chapitre}%
  \renewcommand{\thecoursechapter}{\arabic{coursechapter}}}
\newcommand{\coursepart}[5]{%
  \gdef\coursepartshort{#2}%
  \gdef\coursepartposition{Étape #4}%
  \gdef\coursepartindex{#5}%
  \gdef\coursepartcolor{#3}%
  \setbeamercolor{frametitle}{fg=white,bg=ink}%
  \part{#1}
  \begin{frame}[plain]
    \centering\vfill
    {\scriptsize\bfseries\color{#3}ÉTAPE #4\par}
    \vspace{2mm}
    {\usebeamerfont{part title}\color{ink}#1\par}
    \vspace{3mm}{\color{#3}\rule{28mm}{1pt}}
    \vfill
  \end{frame}}
\newcommand{\coursechapter}[2]{%
  \refstepcounter{coursechapter}%
  \gdef\coursechapterindicator{\coursechapterkind\space\thecoursechapter}%
  \label{#2}%
  \section{#1}
  \begin{frame}[plain]
    \vfill
    {\large\sffamily\bfseries\color{\coursepartcolor}%
      \coursechapterkind\space\thecoursechapter\par}
    \vspace{3mm}
    {\usebeamerfont{section title}\color{ink}#1\par}
    \vspace{4mm}{\color{\coursepartcolor}\rule{24mm}{.9pt}}
    \vfill
  \end{frame}}
\renewcommand{\thefigure}{\arabic{figure}}
\renewcommand{\thetable}{\arabic{table}}

% Ouverture pédagogique d'un chapitre : objectifs scannables, métadonnées après.
\NewDocumentCommand{\orientobjective}{m}{\item #1}
\NewDocumentCommand{\orientobjectives}{>{\SplitList{;}}m}{%
  \begin{itemize}
  \setlength{\itemsep}{.45em}
  \ProcessList{#1}{\orientobjective}
  \end{itemize}}
\newcommand{\orientmeta}[2]{%
  \par\addvspace{3pt}\noindent
  \hangindent=26mm\hangafter=1
  \makebox[26mm][l]{\small\bfseries\color{accent!78!ink}#1}%
  {\small\color{ink!68}#2}\par}
\newcommand{\orientcontent}[3]{%
  \orientobjectives{#1}
  \vspace{.35em}
  \orientmeta{Prérequis}{#2}
  \orientmeta{Parcours}{#3}}

% Encadrés transformés en respirations typographiques qui peuvent se poursuivre
% sur la vue suivante. Aucun grand rectangle ne retient artificiellement le texte.
\newenvironment{coursecallout}[2]{%
  \par\addvspace{.55em}\begingroup
  \leftskip=4mm\rightskip=2mm
  \noindent{\small\bfseries\color{#1}#2}\par\smallskip
  \color{ink!94}\ignorespaces}
  {\par\endgroup\addvspace{.55em}}
\newenvironment{resultat}[1][]{\begin{coursecallout}{accent!82!ink}{#1}}
  {\end{coursecallout}}
\newenvironment{vigilance}[1][]{\begin{coursecallout}{warning!82!ink}{#1}}
  {\end{coursecallout}}
\RenewDocumentEnvironment{lecture}{}{\begin{coursecallout}{ink!68}{Lecture}}
  {\end{coursecallout}}
\newenvironment{chiffres}[1][]{\begin{coursecallout}{ink!58}{#1}\small}
  {\end{coursecallout}}
\newenvironment{releve}{%
  \par\begingroup\small\color{ink!68}\leftskip=5mm\rightskip=2mm}
  {\par\endgroup}
\newenvironment{rappel}[1][]{\begin{coursecallout}{rulecolor!82!ink}{#1}}
  {\end{coursecallout}}
\newenvironment{approfondir}[1][]{\begin{coursecallout}{ink!62}{#1}}
  {\end{coursecallout}}

\NewDocumentEnvironment{panelenumerate}{O{accent}}
  {\begin{enumerate}\setlength{\itemsep}{.45em}}
  {\end{enumerate}}
\NewDocumentCommand{\panelitem}{O{accent}m}{\item #2}
\NewDocumentEnvironment{panelitemize}{O{accent}}
  {\begin{itemize}\setlength{\itemsep}{.45em}}
  {\end{itemize}}
\NewDocumentCommand{\panelbullet}{O{accent}m}{\item #2}
\NewDocumentCommand{\paneltablerow}{O{accent}mmm}{%
  \cellcolor{#1!10}{\bfseries\color{#1!78!ink}#2} &
  \cellcolor{gridgray!6}{\color{ink!94}#3} &
  \cellcolor{gridgray!6}{\color{ink!65}#4} \\}

\newcounter{exo}
\renewcommand{\theexo}{\arabic{exo}}
\NewDocumentEnvironment{exercice}{O{}}{%
  \refstepcounter{exo}\begin{coursecallout}{rulecolor!82!ink}%
    {Exercice~\theexo\space— #1}}
  {\end{coursecallout}}
\newcounter{tache}
\NewDocumentEnvironment{tache}{O{}}{%
  \refstepcounter{tache}\begin{coursecallout}{accent!82!ink}%
    {Étape~\thetache\space— #1}}
  {\end{coursecallout}}
\newcommand{\tlab}[1]{\par\addvspace{3pt}%
  {\scriptsize\bfseries\color{accent!80!ink}#1}\quad\ignorespaces}
\newcommand{\corrige}[1]{\par\addvspace{3pt}%
  {\scriptsize\color{ink!55}Corrigé : \annexeref{ann:corriges},
   \texttt{#1}.}}

\NewDocumentEnvironment{graphique}{m}{}{}
\newcommand{\codefile}[2][firstline=1]{%
  \par\lstinputlisting[style=code,#1]{#2}\par}
\lstnewenvironment{console}
  {\lstset{language=,basicstyle=\ttfamily\scriptsize,
    columns=fullflexible,keepspaces=true}}
  {}

\newcommand{\Gryf}{\textcolor{gryff}{Gryffindor}\xspace}
\newcommand{\Huff}{\textcolor{huffle}{Hufflepuff}\xspace}
\newcommand{\Rav}{\textcolor{raven}{Ravenclaw}\xspace}
\newcommand{\Slyth}{\textcolor{slyth}{Slytherin}\xspace}
\newcommand{\Herb}{\textcolor{herbo}{Herbology}\xspace}
\newcommand{\Runes}{\textcolor{runes}{Ancient Runes}\xspace}
\newcommand{\ann}[3]{%
  \underbrace{\textcolor{#1}{#2}}_{\textcolor{#1}{\substack{#3}}}}
\newcommand{\annup}[3]{%
  \overbrace{\textcolor{#1}{#2}}^{\textcolor{#1}{\substack{#3}}}}
\newcommand{\annl}[1]{\text{\scriptsize #1}}
\newcommand{\R}{\mathbb{R}}
\newcommand{\Pp}{\mathbb{P}}
\newcommand{\x}{\mathbf{x}}
\newcommand{\w}{\mathbf{w}}
\newcommand{\X}{\mathbf{X}}
\newcommand{\y}{\mathbf{y}}
% Les compléments ajoutés tardivement au cours long ne sont pas projetés.
% Les mentions résiduelles restent lisibles sans produire de renvoi cassé.
\newcommand{\annexeref}[1]{annexe du cours long}
\newcommand{\slidebreak}{\framebreak}
\newcommand{\Mzero}{\ensuremath{M_0}\xspace}
\newcommand{\Mquatre}{\ensuremath{M_{400}}\xspace}
\newcommand{\Mstop}{\ensuremath{M_{\text{arr}}}\xspace}
\newcommand{\Movr}{\ensuremath{M_{4}}\xspace}

\hypersetup{
  colorlinks=true,
  linkcolor=accent!85!ink,
  citecolor=accent,
  urlcolor=accent,
  pdftitle={Régression logistique --- diaporama pédagogique},
  pdfauthor={Sylvain Maitre},
  pdflang={fr-FR}}

\begin{document}

\begin{frame}[plain]
  \titlepage
\end{frame}

\begin{frame}{Parcours causal}
  \centering
  {\large\bfseries\color{ink}
    Problème \(\longrightarrow\) Données \(\longrightarrow\) Score
    \(\longrightarrow\) Probabilité\par}
  \vspace{6mm}
  {\large\bfseries\color{ink}
    Perte \(\longrightarrow\) Optimisation
    \(\longrightarrow\) Évaluation\par}
  \vspace{9mm}
  {\small\color{ink!68}
    Chaque notion est introduite, justifiée, programmée puis utilisée avant
    de passer à la suivante.\par}
\end{frame}

% Fichier généré par scripts/gen_slides_source.py.
% Les nombres viennent exclusivement de data/micro_ecoles.csv.
% Fichier genere par scripts/build_micro_case.py ; ne pas modifier.
% Les donnees brutes, l'imputation et les calculs sont controles par le generateur.

\newcommand{\MicroGryffondorLabel}{Gryffondor}
\newcommand{\MicroSerpentardLabel}{Serpentard}
\newcommand{\MicroNTrain}{12}
\newcommand{\MicroNTest}{8}
\newcommand{\MicroMedianPotions}{10}
\newcommand{\MicroMedianVol}{50}
\newcommand{\MicroMeanPotions}{10}
\newcommand{\MicroMeanVol}{50}
\newcommand{\MicroSdPotions}{4}
\newcommand{\MicroSdVol}{12}
\newcommand{\MicroTrainCorrelation}{-\frac{1}{8}}
\newcommand{\MicroExactLogOdds}{\ln 5}
\newcommand{\MicroExactProbabilityHigh}{\frac{5}{6}}
\newcommand{\MicroExactProbabilityLow}{\frac{1}{6}}
\newcommand{\MicroW}{\left(0\,;\,-\frac{2\ln 5}{3}\,;\,+\frac{2\ln 5}{3}\right)}
\newcommand{\MicroWDecimal}{\left(+0{,}000000\,;\,-1{,}072959\,;\,+1{,}072959\right)}
\newcommand{\MicroRawBeta}{\left(-\frac{10\ln 5}{9}\,;\,-\frac{\ln 5}{6}\,;\,+\frac{\ln 5}{18}\right)}
\newcommand{\MicroRawBetaDecimal}{\left(-1{,}788264\,;\,-0{,}268240\,;\,+0{,}089413\right)}
\newcommand{\MicroTrainLoss}{0{,}450561}
\newcommand{\MicroTestLoss}{0{,}457133}
\newcommand{\MicroGDIterations}{152}
\newcommand{\MicroGDLoss}{0{,}450561}
\newcommand{\MicroGDW}{\left(+0{,}000000\,;\,-1{,}072959\,;\,+1{,}072959\right)}
\newcommand{\MicroGradientMaxError}{0{,}00000000001}
\newcommand{\MicroThresholdLow}{0{,}50}
\newcommand{\MicroThresholdHigh}{0{,}65}
\newcommand{\MicroTrainAccuracy}{0{,}8333}
\newcommand{\MicroTestAccuracy}{0{,}8750}
\newcommand{\MicroTestAccuracyTauHigh}{0{,}8750}
\newcommand{\MicroTestPrecision}{0{,}8000}
\newcommand{\MicroTestRecall}{1{,}0000}
\newcommand{\MicroTestSpecificity}{0{,}7500}
\newcommand{\MicroTestFScore}{0{,}8889}
\newcommand{\MicroTestWilsonLow}{0{,}529112}
\newcommand{\MicroTestWilsonHigh}{0{,}977583}
\newcommand{\MicroTrainConfusion}{VP $=5$, FN $=1$, FP $=1$, VN $=5$}
\newcommand{\MicroTestConfusion}{VP $=4$, FN $=0$, FP $=1$, VN $=3$}
\newcommand{\MicroTestConfusionTauHigh}{VP $=3$, FN $=1$, FP $=0$, VN $=4$}
\newcommand{\MicroTrainRows}{
E01 & Harry Potter & Gryffondor & $3$ & $47$ \\
E02 & Drago Malefoy & Serpentard & $9$ & $29$ \\
E03 & Hermione Granger & Gryffondor & $5$ & $53$ \\
E04 & Pansy Parkinson & Serpentard & $11$ & $35$ \\
E05 & Ron Weasley & Gryffondor & $6$ & $56$ \\
E06 & Blaise Zabini & Serpentard & $12$ & $38$ \\
E07 & Neville Londubat & Gryffondor & $14$ & $44$ \\
E08 & Vincent Crabbe & Serpentard & $8$ & $62$ \\
E09 & Ginny Weasley & Gryffondor & $9$ & $65$ \\
E10 & Gregory Goyle & Serpentard & $15$ & $47$ \\
E11 & Dean Thomas & Gryffondor & $11$ & $71$ \\
E12 & Théodore Nott & Serpentard & $17$ & $53$ \\}
\newcommand{\MicroTestRawRows}{
E13 & Seamus Finnigan & Gryffondor & \textcolor{warning}{---} & $60$ \\
E14 & Daphné Greengrass & Serpentard & $16$ & \textcolor{warning}{---} \\
E15 & Parvati Patil & Gryffondor & $7$ & $50$ \\
E16 & Millicent Bulstrode & Serpentard & $7$ & $35$ \\
E17 & Lavande Brown & Gryffondor & $13$ & $65$ \\
E18 & Marcus Flint & Serpentard & $10$ & $55$ \\
E19 & Colin Crivey & Gryffondor & $4$ & $41$ \\
E20 & Adrian Pucey & Serpentard & $13$ & $55$ \\}
\newcommand{\MicroTestPreparedRows}{
E13 & Seamus Finnigan & Gryffondor & $10$ & $60$ \\
E14 & Daphné Greengrass & Serpentard & $16$ & $50$ \\
E15 & Parvati Patil & Gryffondor & $7$ & $50$ \\
E16 & Millicent Bulstrode & Serpentard & $7$ & $35$ \\
E17 & Lavande Brown & Gryffondor & $13$ & $65$ \\
E18 & Marcus Flint & Serpentard & $10$ & $55$ \\
E19 & Colin Crivey & Gryffondor & $4$ & $41$ \\
E20 & Adrian Pucey & Serpentard & $13$ & $55$ \\}
\newcommand{\MicroAllGryffondorCoords}{(3,47) (5,53) (6,56) (14,44) (9,65) (11,71) (10,60) (7,50) (13,65) (4,41)}
\newcommand{\MicroAllSerpentardCoords}{(9,29) (11,35) (12,38) (8,62) (15,47) (17,53) (16,50) (7,35) (10,55) (13,55)}
\newcommand{\MicroTrainGryffondorCoords}{(3,47) (5,53) (6,56) (14,44) (9,65) (11,71)}
\newcommand{\MicroTrainSerpentardCoords}{(9,29) (11,35) (12,38) (8,62) (15,47) (17,53)}
\newcommand{\MicroTestGryffondorCoords}{(10,60) (7,50) (13,65) (4,41)}
\newcommand{\MicroTestSerpentardCoords}{(16,50) (7,35) (10,55) (13,55)}
\newcommand{\MicroTrainGryffondorStdCoords}{(-1.75,-0.25) (-1.25,0.25) (-1,0.5) (1,-0.5) (-0.25,1.25) (0.25,1.75)}
\newcommand{\MicroTrainSerpentardStdCoords}{(-0.25,-1.75) (0.25,-1.25) (0.5,-1) (-0.5,1) (1.25,-0.25) (1.75,0.25)}
\newcommand{\MicroMissingCoords}{(10,60) (16,50)}
\newcommand{\MicroAllNameLabels}{
\node[font=\tiny\bfseries,text=coursescore!88!ink,fill=white,fill opacity=.82,text opacity=1,inner sep=.6pt,xshift=-9pt,yshift=7pt] at (axis cs:3,47) {Harry};
\node[font=\tiny\bfseries,text=warning!88!ink,fill=white,fill opacity=.82,text opacity=1,inner sep=.6pt,xshift=-6pt,yshift=-8pt] at (axis cs:9,29) {Drago};
\node[font=\tiny\bfseries,text=coursescore!88!ink,fill=white,fill opacity=.82,text opacity=1,inner sep=.6pt,xshift=-12pt,yshift=8pt] at (axis cs:5,53) {Hermione};
\node[font=\tiny\bfseries,text=warning!88!ink,fill=white,fill opacity=.82,text opacity=1,inner sep=.6pt,xshift=-9pt,yshift=-8pt] at (axis cs:11,35) {Pansy};
\node[font=\tiny\bfseries,text=coursescore!88!ink,fill=white,fill opacity=.82,text opacity=1,inner sep=.6pt,xshift=10pt,yshift=9pt] at (axis cs:6,56) {Ron};
\node[font=\tiny\bfseries,text=warning!88!ink,fill=white,fill opacity=.82,text opacity=1,inner sep=.6pt,xshift=11pt,yshift=-7pt] at (axis cs:12,38) {Blaise};
\node[font=\tiny\bfseries,text=coursescore!88!ink,fill=white,fill opacity=.82,text opacity=1,inner sep=.6pt,xshift=-11pt,yshift=-9pt] at (axis cs:14,44) {Neville};
\node[font=\tiny\bfseries,text=warning!88!ink,fill=white,fill opacity=.82,text opacity=1,inner sep=.6pt,xshift=-14pt,yshift=10pt] at (axis cs:8,62) {Vincent};
\node[font=\tiny\bfseries,text=coursescore!88!ink,fill=white,fill opacity=.82,text opacity=1,inner sep=.6pt,xshift=13pt,yshift=11pt] at (axis cs:9,65) {Ginny};
\node[font=\tiny\bfseries,text=warning!88!ink,fill=white,fill opacity=.82,text opacity=1,inner sep=.6pt,xshift=-11pt,yshift=10pt] at (axis cs:15,47) {Gregory};
\node[font=\tiny\bfseries,text=coursescore!88!ink,fill=white,fill opacity=.82,text opacity=1,inner sep=.6pt,xshift=0pt,yshift=10pt] at (axis cs:11,71) {Dean};
\node[font=\tiny\bfseries,text=warning!88!ink,fill=white,fill opacity=.82,text opacity=1,inner sep=.6pt,xshift=13pt,yshift=10pt] at (axis cs:17,53) {Théodore};
\node[font=\tiny\bfseries,text=coursescore!88!ink,fill=white,fill opacity=.82,text opacity=1,inner sep=.6pt,xshift=16pt,yshift=-4pt] at (axis cs:10,60) {Seamus};
\node[font=\tiny\bfseries,text=warning!88!ink,fill=white,fill opacity=.82,text opacity=1,inner sep=.6pt,xshift=14pt,yshift=-9pt] at (axis cs:16,50) {Daphné};
\node[font=\tiny\bfseries,text=coursescore!88!ink,fill=white,fill opacity=.82,text opacity=1,inner sep=.6pt,xshift=-11pt,yshift=-9pt] at (axis cs:7,50) {Parvati};
\node[font=\tiny\bfseries,text=warning!88!ink,fill=white,fill opacity=.82,text opacity=1,inner sep=.6pt,xshift=-11pt,yshift=-9pt] at (axis cs:7,35) {Millicent};
\node[font=\tiny\bfseries,text=coursescore!88!ink,fill=white,fill opacity=.82,text opacity=1,inner sep=.6pt,xshift=12pt,yshift=10pt] at (axis cs:13,65) {Lavande};
\node[font=\tiny\bfseries,text=warning!88!ink,fill=white,fill opacity=.82,text opacity=1,inner sep=.6pt,xshift=-17pt,yshift=-9pt] at (axis cs:10,55) {Marcus};
\node[font=\tiny\bfseries,text=coursescore!88!ink,fill=white,fill opacity=.82,text opacity=1,inner sep=.6pt,xshift=-10pt,yshift=-9pt] at (axis cs:4,41) {Colin};
\node[font=\tiny\bfseries,text=warning!88!ink,fill=white,fill opacity=.82,text opacity=1,inner sep=.6pt,xshift=14pt,yshift=-8pt] at (axis cs:13,55) {Adrian};}
\newcommand{\MicroTrainAtypicalCoords}{(14,44) (8,62)}
\newcommand{\MicroTestAtypicalCoords}{(10,55)}
\newcommand{\MicroGDTableRows}{
$0$ & $0{,}693147$ & $+0{,}000000$ & $+0{,}000000$ & $+0{,}000000$ & $0{,}353553$ \\
$1$ & $0{,}585623$ & $+0{,}000000$ & $-0{,}250000$ & $+0{,}250000$ & $0{,}255266$ \\
$2$ & $0{,}529140$ & $+0{,}000000$ & $-0{,}430500$ & $+0{,}430500$ & $0{,}188035$ \\
$5$ & $0{,}469709$ & $+0{,}000000$ & $-0{,}741135$ & $+0{,}741135$ & $0{,}085794$ \\
$10$ & $0{,}453187$ & $+0{,}000000$ & $-0{,}946007$ & $+0{,}946007$ & $0{,}029858$ \\
$20$ & $0{,}450636$ & $+0{,}000000$ & $-1{,}051109$ & $+1{,}051109$ & $0{,}004881$ \\
$50$ & $0{,}450561$ & $+0{,}000000$ & $-1{,}072827$ & $+1{,}072827$ & $2{,}915\times10^{-5}$ \\
$100$ & $0{,}450561$ & $+0{,}000000$ & $-1{,}072959$ & $+1{,}072959$ & $5{,}961\times10^{-9}$ \\
$152$ & $0{,}450561$ & $+0{,}000000$ & $-1{,}072959$ & $+1{,}072959$ & $8{,}678\times10^{-13}$ \\}
\newcommand{\MicroGradientCheckRows}{
$0$ & $+0{,}046541807$ & $+0{,}046541807$ & $0{,}00000000000$ \\
$1$ & $+0{,}144621051$ & $+0{,}144621051$ & $0{,}00000000000$ \\
$2$ & $-0{,}164881324$ & $-0{,}164881324$ & $0{,}00000000001$ \\}
\newcommand{\MicroMetricsRows}{
apprentissage & $0{,}50$ & $5$ & $1$ & $1$ & $5$ & $0{,}8333$ & $0{,}8333$ & $0{,}8333$ & $0{,}8333$ & $0{,}8333$ & $0{,}8333$ \\
apprentissage & $0{,}65$ & $5$ & $1$ & $1$ & $5$ & $0{,}8333$ & $0{,}8333$ & $0{,}8333$ & $0{,}8333$ & $0{,}8333$ & $0{,}8333$ \\
test & $0{,}50$ & $4$ & $0$ & $1$ & $3$ & $0{,}8750$ & $0{,}8000$ & $1{,}0000$ & $0{,}7500$ & $0{,}8889$ & $0{,}8750$ \\
test & $0{,}65$ & $3$ & $1$ & $0$ & $4$ & $0{,}8750$ & $1{,}0000$ & $0{,}7500$ & $1{,}0000$ & $0{,}8571$ & $0{,}8750$ \\}
\newcommand{\MicroTestSummaryRows}{
$0{,}50$ & $0{,}8750$ & $0{,}8000$ & $1{,}0000$ & $0{,}7500$ & $0{,}8889$ \\
$0{,}65$ & $0{,}8750$ & $1{,}0000$ & $0{,}7500$ & $1{,}0000$ & $0{,}8571$ \\}
\newcommand{\MicroTrainCalcRows}{
E01 & Harry & Gryff. & $+1{,}609438$ & $0{,}833333$ & Gryff. & Gryff. \\
E02 & Drago & Serp. & $-1{,}609438$ & $0{,}166667$ & Serp. & Serp. \\
E03 & Hermione & Gryff. & $+1{,}609438$ & $0{,}833333$ & Gryff. & Gryff. \\
E04 & Pansy & Serp. & $-1{,}609438$ & $0{,}166667$ & Serp. & Serp. \\
E05 & Ron & Gryff. & $+1{,}609438$ & $0{,}833333$ & Gryff. & Gryff. \\
E06 & Blaise & Serp. & $-1{,}609438$ & $0{,}166667$ & Serp. & Serp. \\
E07 & Neville & Gryff. & $-1{,}609438$ & $0{,}166667$ & Serp. & Serp. \\
E08 & Vincent & Serp. & $+1{,}609438$ & $0{,}833333$ & Gryff. & Gryff. \\
E09 & Ginny & Gryff. & $+1{,}609438$ & $0{,}833333$ & Gryff. & Gryff. \\
E10 & Gregory & Serp. & $-1{,}609438$ & $0{,}166667$ & Serp. & Serp. \\
E11 & Dean & Gryff. & $+1{,}609438$ & $0{,}833333$ & Gryff. & Gryff. \\
E12 & Théodore & Serp. & $-1{,}609438$ & $0{,}166667$ & Serp. & Serp. \\}
\newcommand{\MicroTestCalcRows}{
E13 & Seamus & Gryff. & $+0{,}894132$ & $0{,}709742$ & Gryff. & Gryff. \\
E14 & Daphné & Serp. & $-1{,}609438$ & $0{,}166667$ & Serp. & Serp. \\
E15 & Parvati & Gryff. & $+0{,}804719$ & $0{,}690983$ & Gryff. & Gryff. \\
E16 & Millicent & Serp. & $-0{,}536479$ & $0{,}369007$ & Serp. & Serp. \\
E17 & Lavande & Gryff. & $+0{,}536479$ & $0{,}630993$ & Gryff. & Serp. \\
E18 & Marcus & Serp. & $+0{,}447066$ & $0{,}609941$ & Gryff. & Serp. \\
E19 & Colin & Gryff. & $+0{,}804719$ & $0{,}690983$ & Gryff. & Gryff. \\
E20 & Adrian & Serp. & $-0{,}357653$ & $0{,}411528$ & Serp. & Serp. \\}
\newcommand{\MicroProbabilityLavande}{0{,}630993}
\newcommand{\MicroProbabilityMarcus}{0{,}609941}

% Contenu éditorial du cas Gryffondor--Serpentard.
% Toutes les valeurs numériques affichées sont fournies par le générateur.

\newcommand{\GryffondorName}{\textcolor{coursescore!88!ink}{\MicroGryffondorLabel}}
\newcommand{\SerpentardName}{\textcolor{warning!88!ink}{\MicroSerpentardLabel}}
\newcommand{\microheader}{\rowcolor{ink!92}\color{white}\bfseries}

\pgfplotsset{
  rawcloud/.style={
    width=.86\textwidth,height=5.65cm,
    xmin=0,xmax=20,ymin=0,ymax=100,
    xtick={0,5,10,15,20},ytick={0,20,40,60,80,100},
    axis lines=box,axis line style={ink!42},
    grid=major,grid style={gridgray!34},
    xlabel={Potions (sur 20)},ylabel={Vol (sur 100)},
    tick label style={font=\scriptsize},label style={font=\scriptsize},
    legend style={at={(.5,1.02)},anchor=south,draw=none,
      font=\scriptsize,legend columns=3,
      /tikz/every even column/.append style={column sep=7pt}}},
  stdcloud/.style={
    width=\textwidth,height=4.8cm,
    xmin=-2,xmax=2,ymin=-2,ymax=2,
    xtick={-2,-1,0,1,2},ytick={-2,-1,0,1,2},
    axis lines=box,axis line style={ink!42},
    grid=major,grid style={gridgray!34},
    xlabel={$x_{\mathrm{pot}}$},ylabel={$x_{\mathrm{vol}}$},
    tick label style={font=\tiny},label style={font=\scriptsize}}}

\courseintro

\begin{frame}[t]{Formulation du problème de classification}
  \vspace{4mm}
  \begin{columns}[T,onlytextwidth]
    \column{.47\textwidth}
      \begin{coursecallout}{coursedata!88!ink}{Observations disponibles}
        Pour chaque élève :
        \begin{itemize}
          \item une note de Potions sur 20 ;
          \item une note de Vol sur 100.
        \end{itemize}
      \end{coursecallout}
    \column{.47\textwidth}
      \begin{coursecallout}{coursedecision!88!ink}{Maison observée}
        Deux maisons sont représentées :
        \begin{itemize}
          \item \GryffondorName ;
          \item \SerpentardName.
        \end{itemize}
      \end{coursecallout}
  \end{columns}
  \vfill
  \centering
  {\large\bfseries Construire une règle qui associe les deux notes à une maison.}
\end{frame}

\begin{frame}[t]{Distribution des notes de Potions et de Vol}
  \centering
  \begin{tikzpicture}
    \begin{axis}[rawcloud,width=.94\textwidth,height=6.25cm]
      \addplot[only marks,mark=square*,mark size=2.5pt,coursescore!88!ink]
        coordinates {\MicroAllGryffondorCoords};
      \addlegendentry{Gryffondor}
      \addplot[only marks,mark=*,mark size=2.5pt,warning!88!ink]
        coordinates {\MicroAllSerpentardCoords};
      \addlegendentry{Serpentard}
      \addplot[only marks,mark=o,mark size=4.4pt,very thick,ink!72]
        coordinates {\MicroMissingCoords};
      \addlegendentry{note imputée}
      \MicroAllNameLabels
    \end{axis}
  \end{tikzpicture}
\end{frame}

\begin{frame}[t]{Chaîne de construction du classifieur}
\centering
\begin{slidevisual}
\begin{tikzpicture}[scale=.97,transform shape,
  flow/.style={-{Latex[length=1.8mm]},line width=.75pt,ink!62},
  card/.style={draw=ink!24,fill=white,rounded corners=1.6pt,line width=.5pt},
  head/.style={font=\fontsize{6.3}{7}\selectfont\bfseries,text=ink!92}]
  \def\cw{3.25}
  \def\ch{2.55}
  \foreach \x/\col/\title in {
    0/coursedata/DONNÉES,
    3.8/coursescore/SCORE LINÉAIRE,
    7.6/courseprobability/SIGMOÏDE,
    11.4/coursedecision/DÉCISION}
  {
    \path[card] (\x,4.15) rectangle ++(\cw,\ch);
    \path[fill=\col!16,rounded corners=1.6pt]
      (\x,6.19) rectangle ++(\cw,.51);
    \node[head,anchor=west] at (\x+.16,6.445) {\title};
  }
  \foreach \x/\col/\title in {
    0/courseupdate/MISE À JOUR,
    3.8/coursegradient/GRADIENT,
    7.6/courseloss/PERTE LOGISTIQUE,
    11.4/courseevaluation/ÉVALUATION}
  {
    \path[card] (\x,.35) rectangle ++(\cw,\ch);
    \path[fill=\col!16,rounded corners=1.6pt]
      (\x,2.39) rectangle ++(\cw,.51);
    \node[head,anchor=west] at (\x+.16,2.645) {\title};
  }

  % Données : deux échelles et une valeur absente.
  \draw[ink!55,line width=.45pt] (.42,4.55) -- (.42,5.93);
  \draw[ink!55,line width=.45pt] (1.25,4.55) -- (1.25,5.93);
  \fill[coursedata!62] (.58,4.55) rectangle (.86,5.25);
  \fill[coursedata!32] (1.41,4.55) rectangle (1.69,5.72);
  \draw[warning,dashed,line width=.7pt] (2.13,4.77) rectangle (2.72,5.35);
  \node[font=\tiny,text=ink!65] at (.72,5.47) {/20};
  \node[font=\tiny,text=ink!65] at (1.55,5.92) {/100};
  \node[font=\tiny,text=warning] at (2.425,5.06) {manq.};

  % Score : nuage et niveau nul oblique.
  \draw[-{Latex[length=1.1mm]},ink!52] (4.19,4.55) -- (6.64,4.55);
  \draw[-{Latex[length=1.1mm]},ink!52] (4.19,4.55) -- (4.19,5.95);
  \draw[coursescore!80!ink,line width=.8pt] (4.45,4.88) -- (6.42,5.65);
  \foreach \p in {(4.55,5.42),(5.06,5.66),(5.55,5.77)}
    \node[fill=coursescore,minimum size=2.5pt,inner sep=0pt] at \p {};
  \foreach \p in {(4.72,4.73),(5.38,4.92),(6.05,5.09)}
    \fill[warning] \p circle (1.25pt);

  % Sigmoïde.
  \begin{scope}[shift={(9.23,4.72)},x=.38cm,y=.9cm]
    \draw[-{Latex[length=1mm]},ink!50] (-3.1,0) -- (3.2,0);
    \draw[-{Latex[length=1mm]},ink!50] (0,-.06) -- (0,1.15);
    \draw[courseprobability!88!ink,line width=.9pt]
      plot[samples=60,domain=-3:3] (\x,{1/(1+exp(-\x))});
    \fill[courseprobability] (0,.5) circle (1.4pt);
  \end{scope}

  % Décision.
  \path[fill=warning!14] (11.82,4.78) rectangle (13.02,5.65);
  \path[fill=coursescore!14] (13.02,4.78) rectangle (14.23,5.65);
  \draw[ink!68,line width=.7pt] (13.02,4.65) -- (13.02,5.8);
  \node[font=\tiny\bfseries,text=warning!88!ink] at (12.42,5.23) {SERP.};
  \node[font=\tiny\bfseries,text=coursescore!88!ink] at (13.62,5.23) {GRYFF.};
  \node[font=\tiny,text=ink!65] at (13.02,4.57) {$\tau$};

  % Mise à jour : trajectoire de trois coefficients.
  \draw[courseupdate!35,line width=.45pt] (1.28,1.22) ellipse (.98cm and .58cm);
  \draw[courseupdate!55,line width=.5pt] (1.28,1.22) ellipse (.61cm and .36cm);
  \fill[ink!72] (2.32,1.95) circle (1.5pt);
  \fill[ink!55] (1.85,1.62) circle (1.4pt);
  \fill[courseupdate] (1.48,1.34) circle (1.5pt);
  \draw[flow,courseupdate!85!ink] (2.32,1.95) -- (1.89,1.65);
  \draw[flow,courseupdate!85!ink] (1.79,1.57) -- (1.52,1.37);

  % Gradient : courbes de niveau et flèche locale.
  \draw[coursegradient!35,line width=.45pt] (5.42,1.62) ellipse (1.1cm and .67cm);
  \draw[coursegradient!55,line width=.5pt] (5.42,1.62) ellipse (.68cm and .4cm);
  \fill[ink!74] (6.18,2.06) circle (1.5pt);
  \draw[-{Latex[length=1.5mm]},coursegradient!88!ink,line width=.9pt]
    (6.18,2.06) -- (5.68,1.76);
  \node[font=\tiny,text=coursegradient!88!ink] at (6.18,1.72) {$-\nabla J$};

  % Perte : deux courbes opposées, sans codage numérique.
  \begin{scope}[shift={(9.22,.66)},x=.36cm,y=.43cm]
    \draw[-{Latex[length=1mm]},ink!50] (-3.1,0) -- (3.2,0);
    \draw[-{Latex[length=1mm]},ink!50] (-3.0,0) -- (-3.0,3.25);
    \draw[coursescore!88!ink,line width=.75pt]
      plot[samples=50,domain=-2.9:2.9] (\x,{ln(1+exp(-\x))});
    \draw[warning!88!ink,line width=.75pt]
      plot[samples=50,domain=-2.9:2.9] (\x,{ln(1+exp(\x))});
    \node[font=\tiny,text=coursescore!88!ink] at (-2.15,2.72) {G};
    \node[font=\tiny,text=warning!88!ink] at (2.15,2.72) {S};
  \end{scope}

  % Évaluation : matrice 2 x 2.
  \foreach \x/\y/\col in {11.95/1.64/rulecolor,12.72/1.64/warning,
                            11.95/.88/warning,12.72/.88/rulecolor}
    \path[fill=\col!18,draw=white] (\x,\y) rectangle ++(.72,.67);
  \node[font=\tiny,text=ink!68] at (12.69,.60)
    {observé / prédit};

  \draw[flow] (3.25,5.43) -- (3.8,5.43);
  \draw[flow] (7.05,5.43) -- (7.6,5.43);
  \draw[flow] (10.85,5.43) -- (11.4,5.43);
  \draw[flow] (9.225,4.15) -- (9.225,2.9);
  \draw[flow] (7.6,1.62) -- (7.05,1.62);
  \draw[flow] (3.8,1.62) -- (3.25,1.62);
  \draw[flow] (1.625,2.9) -- (1.625,3.45) -- (5.425,3.45) -- (5.425,4.15);
  \draw[flow] (13.025,4.15) -- (13.025,2.9);
\end{tikzpicture}
\end{slidevisual}
\end{frame}

% ---------------------------------------------------------------------------
\coursepart{Définir les données et le protocole}{Données}{coursedata}

\begin{frame}[t]{Séparation entre estimation et évaluation}
  \vspace{5mm}
  \centering
  \begin{tikzpicture}[node distance=9mm and 12mm,
    box/.style={draw=ink!28,rounded corners=2pt,minimum height=15mm,
      minimum width=31mm,align=center,font=\small,inner xsep=4mm},
    arr/.style={-{Latex[length=2mm]},thick,ink!58}]
    \node[box,fill=coursedata!13] (all) {20 élèves};
    \node[box,fill=coursedata!22,right=of all,yshift=10mm] (train)
      {12 élèves\\estimation};
    \node[box,fill=courseevaluation!18,right=of all,yshift=-10mm] (test)
      {8 élèves\\évaluation};
    \node[box,fill=courseupdate!15,right=of train] (fit)
      {médianes, $\mu$, $s$\\et coefficients};
    \node[box,fill=courseevaluation!24,right=of test] (eval)
      {mesure\\des erreurs};
    \draw[arr] (all) -- (train);
    \draw[arr] (all) -- (test);
    \draw[arr] (train) -- (fit);
    \draw[arr] (fit.south) |- (eval.north);
    \draw[arr] (test) -- (eval);
  \end{tikzpicture}
  \vfill
  Le groupe d'estimation constitue le \emph{jeu d'apprentissage} : lui seul fixe
  les nombres du modèle. Le groupe d'évaluation les reçoit sans les modifier.
\end{frame}

\begin{frame}[t]{Groupe d'estimation : douze élèves}
  \centering
  \scriptsize
  \renewcommand{\arraystretch}{1.04}
  \begin{tabularx}{\textwidth}{c >{\raggedright\arraybackslash}X l c c}
    \microheader ID & \color{white}\bfseries Élève &
      \color{white}\bfseries Maison & \color{white}\bfseries Potions &
      \color{white}\bfseries Vol \\
    \MicroTrainRows
  \end{tabularx}
\end{frame}

\begin{frame}[t]{Groupe d'évaluation : huit élèves}
  \centering
  \small
  \renewcommand{\arraystretch}{1.04}
  \begin{tabularx}{\textwidth}{c >{\raggedright\arraybackslash}X l c c}
    \microheader ID & \color{white}\bfseries Élève &
      \color{white}\bfseries Maison & \color{white}\bfseries Potions &
      \color{white}\bfseries Vol \\
    \MicroTestRawRows
  \end{tabularx}
  \vfill
  \begin{vigilance}[Deux notes sont manquantes]
    Potions manque pour Seamus ; Vol manque pour Daphné. Une case vide ne peut
    pas être standardisée : elle doit d'abord recevoir une valeur numérique.
  \end{vigilance}
\end{frame}

\begin{frame}[t]{Compléter, puis standardiser une note manquante}
  \vspace{3mm}
  \begin{columns}[T,onlytextwidth]
    \column{.47\textwidth}
      \begin{coursecallout}{coursedata!88!ink}{Valeurs de remplacement}
        \[
          \widetilde P=\MicroMedianPotions,
          \qquad \widetilde V=\MicroMedianVol.
        \]
        Ces médianes sont calculées sur les douze élèves du groupe d'estimation.
      \end{coursecallout}
    \column{.47\textwidth}
      \begin{coursecallout}{courseevaluation!88!ink}{Deux substitutions}
        \[
          \begin{aligned}
            P_{\text{Seamus}}&\leftarrow\MicroMedianPotions
              &\Longrightarrow x_{\mathrm{pot}}&=0,\\
            V_{\text{Daphné}}&\leftarrow\MicroMedianVol
              &\Longrightarrow x_{\mathrm{vol}}&=0.
          \end{aligned}
        \]
        Ici, médiane et moyenne coïncident dans chaque matière.
      \end{coursecallout}
  \end{columns}
  \vfill
  \centering
  \begin{adjustbox}{max width=.96\textwidth,center}
  \begin{tikzpicture}[node distance=7mm,
    box/.style={draw=ink!25,rounded corners=2pt,minimum height=11mm,
      align=center,font=\small,inner xsep=4mm},
    arr/.style={-{Latex[length=1.8mm]},thick,ink!58}]
    \node[box,fill=warning!12] (raw) {note manquante};
    \node[box,fill=coursedata!16,right=of raw] (rule) {médiane\\de la matière};
    \node[box,fill=rulecolor!16,right=of rule] (ready) {valeur\\numérique $v$};
    \node[box,fill=courseprobability!14,right=of ready] (standard)
      {$x=(v-\mu)/s$\\note standardisée};
    \draw[arr] (raw) -- (rule);
    \draw[arr] (rule) -- (ready);
    \draw[arr] (ready) -- (standard);
  \end{tikzpicture}
  \end{adjustbox}
\end{frame}

\begin{frame}[t]{Deux matières, deux échelles de mesure}
  \vspace{2mm}
  \begin{columns}[T,onlytextwidth]
    \column{.46\textwidth}
      \centering
      {\bfseries\color{coursedata!88!ink} Potions}
      \begin{tikzpicture}[y=.045cm]
        \draw[-{Latex[length=1.7mm]},ink!58] (0,0) -- (0,23);
        \foreach \v in {0,5,10,15,20}
          \draw[ink!55] (-.12,\v) -- (.12,\v)
            node[right=2mm,font=\scriptsize] {\v};
        \fill[coursedata!68] (-.35,3) rectangle (.35,17);
      \end{tikzpicture}
      \par notée sur 20
    \column{.46\textwidth}
      \centering
      {\bfseries\color{courseprobability!88!ink} Vol}
      \begin{tikzpicture}[y=.009cm]
        \draw[-{Latex[length=1.7mm]},ink!58] (0,0) -- (0,113);
        \foreach \v in {0,25,50,75,100}
          \draw[ink!55] (-.12,\v) -- (.12,\v)
            node[right=2mm,font=\scriptsize] {\v};
        \fill[courseprobability!58] (-.35,29) rectangle (.35,71);
      \end{tikzpicture}
      \par notée sur 100
  \end{columns}
  \vfill
  Un point représente 5\% de l'échelle en Potions, contre 1\% en Vol.
\end{frame}

\begin{frame}[t]{Centrage-réduction des deux notes}
  \vspace{2mm}
  Les paramètres de standardisation sont calculés sur le groupe d'estimation,
  puis conservés :
  \begin{columns}[T,onlytextwidth]
    \column{.47\textwidth}
      \begin{coursecallout}{coursedata!88!ink}{Potions}
        \[
          \mu_{\mathrm{pot}}=\MicroMeanPotions,
          \quad s_{\mathrm{pot}}=\MicroSdPotions,
        \]
        \[
          x_{\mathrm{pot}}=
          \frac{P-\MicroMeanPotions}{\MicroSdPotions}.
        \]
      \end{coursecallout}
    \column{.47\textwidth}
      \begin{coursecallout}{courseprobability!88!ink}{Vol}
        \[
          \mu_{\mathrm{vol}}=\MicroMeanVol,
          \quad s_{\mathrm{vol}}=\MicroSdVol,
        \]
        \[
          x_{\mathrm{vol}}=
          \frac{V-\MicroMeanVol}{\MicroSdVol}.
        \]
      \end{coursecallout}
  \end{columns}
  \vfill
  Pour Harry : \(x_{\mathrm{pot}}=-1{,}75\) et
  \(x_{\mathrm{vol}}=-0{,}25\). Les deux coordonnées sont désormais sans unité.
\end{frame}

\begin{frame}[t]{Effet géométrique de la standardisation}
  \centering
  \begin{columns}[T,onlytextwidth]
    \column{.49\textwidth}
      \centering\small Notes brutes
      \begin{tikzpicture}
        \begin{axis}[width=\textwidth,height=4.8cm,
          xmin=0,xmax=20,ymin=0,ymax=100,
          xtick={0,10,20},ytick={0,50,100},
          grid=major,grid style={gridgray!34},axis lines=box,
          xlabel={Potions},ylabel={Vol},
          tick label style={font=\tiny},label style={font=\scriptsize}]
          \addplot[only marks,mark=square*,mark size=2.2pt,coursescore]
            coordinates {\MicroTrainGryffondorCoords};
          \addplot[only marks,mark=*,mark size=2.2pt,warning]
            coordinates {\MicroTrainSerpentardCoords};
        \end{axis}
      \end{tikzpicture}
    \column{.49\textwidth}
      \centering\small Coordonnées standardisées
      \begin{tikzpicture}
        \begin{axis}[stdcloud]
          \addplot[only marks,mark=square*,mark size=2.2pt,coursescore]
            coordinates {\MicroTrainGryffondorStdCoords};
          \addplot[only marks,mark=*,mark size=2.2pt,warning]
            coordinates {\MicroTrainSerpentardStdCoords};
        \end{axis}
      \end{tikzpicture}
  \end{columns}
  \vfill
  Le centrage-réduction change les unités, pas l'identité des élèves ni leur maison.
\end{frame}

% ---------------------------------------------------------------------------
\coursepart{Construire le score linéaire}{Score linéaire}{coursescore}

\begin{frame}[t]{Définition du score linéaire}
  \vspace{5mm}
  \[
    \boxed{z_i=w_0+w_{\mathrm{pot}}x_{i,\mathrm{pot}}
                    +w_{\mathrm{vol}}x_{i,\mathrm{vol}}}
  \]
  \vspace{4mm}
  \begin{columns}[T,onlytextwidth]
    \column{.47\textwidth}
      \begin{coursecallout}{coursescore!88!ink}{Coefficients}
        Le signe indique le sens de l'association ; la valeur absolue en mesure
        l'intensité sur l'échelle standardisée.
      \end{coursecallout}
    \column{.47\textwidth}
      \begin{coursecallout}{coursegradient!88!ink}{État du modèle}
        Les coefficients sont encore inconnus. Leur estimation viendra après
        la définition de la perte et du gradient.
      \end{coursecallout}
  \end{columns}
\end{frame}

\begin{frame}[t]{Contributions des deux matières au score}
  \centering
  Prenons provisoirement \(w^{\mathrm{essai}}=(0,-1,+1)\), donc
  \(z=x_{\mathrm{vol}}-x_{\mathrm{pot}}\).
  \vspace{3mm}
  \renewcommand{\arraystretch}{1.25}
  \begin{tabular}{lrrrrr}
    \toprule
    élève & $x_{\mathrm{pot}}$ & $x_{\mathrm{vol}}$ &
      $-x_{\mathrm{pot}}$ & $+x_{\mathrm{vol}}$ & $z$\\
    \midrule
    Harry   & $-1{,}75$ & $-0{,}25$ & $+1{,}75$ & $-0{,}25$ & $+1{,}50$\\
    Drago   & $-0{,}25$ & $-1{,}75$ & $+0{,}25$ & $-1{,}75$ & $-1{,}50$\\
    Neville & $+1{,}00$ & $-0{,}50$ & $-1{,}00$ & $-0{,}50$ & $-1{,}50$\\
    \bottomrule
  \end{tabular}
  \vfill
  Le score agrège les contributions ; il ne constitue encore ni une
  probabilité ni une classe prédite.
\end{frame}

\begin{frame}[t]{Le niveau de score nul définit une droite}
  \centering
  \begin{tikzpicture}
    \begin{axis}[rawcloud,legend columns=2]
      \addplot[only marks,mark=square*,mark size=2.6pt,coursescore!88!ink]
        coordinates {\MicroTrainGryffondorCoords};
      \addlegendentry{Gryffondor}
      \addplot[only marks,mark=*,mark size=2.6pt,warning!88!ink]
        coordinates {\MicroTrainSerpentardCoords};
      \addlegendentry{Serpentard}
      \addplot[ink!78,very thick,domain=0:20,samples=2] {3*x+20};
      \node[font=\scriptsize,text=coursescore!88!ink,anchor=south]
        at (axis cs:15,70) {$z>0$};
      \node[font=\scriptsize,text=warning!88!ink,anchor=north]
        at (axis cs:5,25) {$z<0$};
      \node[font=\scriptsize,fill=white,inner sep=1pt,text=ink!72]
        at (axis cs:15.5,64) {$V=3P+20$};
    \end{axis}
  \end{tikzpicture}
\end{frame}

\begin{frame}[t]{Écriture matricielle des scores}
  Pour les douze observations du groupe d'estimation :
  \[
    X=
    \begin{pmatrix}
      1 & x_{1,\mathrm{pot}} & x_{1,\mathrm{vol}}\\
      \vdots & \vdots & \vdots\\
      1 & x_{12,\mathrm{pot}} & x_{12,\mathrm{vol}}
    \end{pmatrix},
    \qquad
    w=\begin{pmatrix}w_0\\w_{\mathrm{pot}}\\w_{\mathrm{vol}}\end{pmatrix}.
  \]
  \[
    \boxed{z=Xw}
  \]
  \vfill
  Une même multiplication produit les douze scores ; chaque ligne de \(X\)
  correspond à un élève.
\end{frame}

% ---------------------------------------------------------------------------
\coursepart{Transformer le score en probabilité}{Sigmoïde}{courseprobability}

\begin{frame}[t]{La sigmoïde associe une probabilité à tout score réel}
  \centering
  \begin{tikzpicture}
    \begin{axis}[
      width=.78\textwidth,height=5.5cm,
      xmin=-5,xmax=5,ymin=-.05,ymax=1.05,
      axis lines=middle,grid=major,grid style={gridgray!35},
      xlabel={$z$},ylabel={$p$},
      xtick={-4,-2,0,2,4},ytick={0,.5,1},
      tick label style={font=\scriptsize},label style={font=\scriptsize}]
      \addplot[courseprobability!88!ink,very thick,domain=-5:5,samples=120]
        {1/(1+exp(-x))};
      \addplot[ink!45,dashed,domain=-5:5,samples=2] {.5};
      \addplot[only marks,mark=*,mark size=2.5pt,coursescore]
        coordinates {(0,.5)};
    \end{axis}
  \end{tikzpicture}
  \[
    \boxed{p_i=\sigma(z_i)=\frac{1}{1+e^{-z_i}}}
  \]
\end{frame}

\begin{frame}[t]{Propriétés de la fonction logistique}
  \vspace{4mm}
  \begin{columns}[T,onlytextwidth]
    \column{.31\textwidth}
      \begin{coursecallout}{courseprobability!88!ink}{Monotonie}
        \[
          z_1<z_2\Rightarrow\sigma(z_1)<\sigma(z_2).
        \]
      \end{coursecallout}
    \column{.31\textwidth}
      \begin{coursecallout}{coursescore!88!ink}{Point central}
        \[
          \sigma(0)=\frac12.
        \]
      \end{coursecallout}
    \column{.31\textwidth}
      \begin{coursecallout}{courseloss!88!ink}{Symétrie}
        \[
          \sigma(-z)=1-\sigma(z).
        \]
      \end{coursecallout}
  \end{columns}
  \vfill
  \centering
  La sigmoïde conserve l'ordre des scores et les transforme en valeurs
  strictement comprises entre \(0\) et \(1\).
\end{frame}

\begin{frame}[t]{Le logit est l'inverse de la sigmoïde}
  \vspace{5mm}
  \begin{columns}[T,onlytextwidth]
    \column{.47\textwidth}
      \begin{coursecallout}{courseprobability!88!ink}{Du score à la probabilité}
        \[
          \frac{p}{1-p}=e^z,
          \qquad p=\frac{e^z}{1+e^z}.
        \]
      \end{coursecallout}
    \column{.47\textwidth}
      \begin{coursecallout}{coursescore!88!ink}{De la probabilité au score}
        \[
          \operatorname{logit}(p)
          =\ln\!\left(\frac{p}{1-p}\right)=z.
        \]
      \end{coursecallout}
  \end{columns}
  \vfill
  \centering
  \begin{tabular}{c|ccccc}
    $z$ & $-\ln5$ & $-1$ & $0$ & $1$ & $+\ln5$\\
    \midrule
    $p$ & $1/6$ & $0{,}269$ & $1/2$ & $0{,}731$ & $5/6$
  \end{tabular}
\end{frame}

\begin{frame}[t]{Probabilités associées à trois scores}
  \centering
  \renewcommand{\arraystretch}{1.35}
  \begin{tabular}{lccc}
    \toprule
    élève & $z$ d'essai & $p=\sigma(z)$ & maison favorisée\\
    \midrule
    Harry   & $+1{,}50$ & $0{,}818$ & Gryffondor\\
    Drago   & $-1{,}50$ & $0{,}182$ & Serpentard\\
    Neville & $-1{,}50$ & $0{,}182$ & Serpentard\\
    \bottomrule
  \end{tabular}
  \vfill
  \begin{coursecallout}{courseprobability!88!ink}{}
    La sigmoïde étant strictement croissante, deux observations de même score
    reçoivent la même probabilité selon le modèle.
  \end{coursecallout}
\end{frame}

% ---------------------------------------------------------------------------
\coursepart{Définir la perte logistique}{Perte logistique}{courseloss}

\begin{frame}[t]{Codage de la maison observée}
  \vspace{5mm}
  \[
    y_i=
    \begin{cases}
      1 & \text{si l'élève appartient à Gryffondor},\\
      0 & \text{si l'élève appartient à Serpentard}.
    \end{cases}
  \]
  \vspace{5mm}
  \begin{resultat}
    Ce codage n'était nécessaire ni pour calculer le score \(z_i\), ni pour le
    transformer en probabilité \(p_i\). Il devient nécessaire pour comparer la
    sortie du modèle à la maison observée.
  \end{resultat}
\end{frame}

\begin{frame}[t]{Perte logistique d'une observation}
  \vspace{3mm}
  \[
    \boxed{\ell(z,y)=-y\ln p-(1-y)\ln(1-p)}
  \]
  \begin{columns}[T,onlytextwidth]
    \column{.47\textwidth}
      \begin{coursecallout}{coursescore!88!ink}{Maison observée : Gryffondor}
        \[
          y=1,\qquad \ell=-\ln p=\ln(1+e^{-z}).
        \]
        La perte décroît lorsque le score augmente.
      \end{coursecallout}
    \column{.47\textwidth}
      \begin{coursecallout}{warning!88!ink}{Maison observée : Serpentard}
        \[
          y=0,\qquad \ell=-\ln(1-p)=\ln(1+e^z).
        \]
        La perte décroît lorsque le score diminue.
      \end{coursecallout}
  \end{columns}
\end{frame}

\begin{frame}[t]{Comportement de la perte selon la maison observée}
  \centering
  \begin{tikzpicture}
    \begin{axis}[width=.74\textwidth,height=5.2cm,xmin=-4,xmax=4,ymin=0,ymax=4.2,
      axis lines=box,grid=major,grid style={gridgray!30},
      xlabel={$z$},ylabel={$\ell$},legend style={draw=none,font=\scriptsize}]
      \addplot[coursescore!88!ink,thick,domain=-4:4,samples=100]
        {ln(1+exp(-x))};
      \addlegendentry{Gryffondor observé}
      \addplot[warning!88!ink,thick,domain=-4:4,samples=100]
        {ln(1+exp(x))};
      \addlegendentry{Serpentard observé}
    \end{axis}
  \end{tikzpicture}
  \[
    \ell(z,y)=\operatorname{softplus}(z)-yz.
  \]
\end{frame}

\begin{frame}[t]{Influence des profils en recouvrement}
  Avec les coefficients d'essai, le signe du score concorde avec la maison
  observée pour dix élèves et s'y oppose pour deux :
  \[
    J=\frac{10\times0{,}201+2\times1{,}701}{12}=0{,}451.
  \]
  \vspace{3mm}
  \begin{columns}[T,onlytextwidth]
    \column{.47\textwidth}
      \begin{coursecallout}{rulecolor!88!ink}{Signe concordant}
        Perte : \(0{,}201\).
      \end{coursecallout}
    \column{.47\textwidth}
      \begin{coursecallout}{warning!88!ink}{Signe discordant}
        Perte : \(1{,}701\), soit environ \(8{,}4\) fois plus.
      \end{coursecallout}
  \end{columns}
  \vfill
  Ces deux observations empêchent une séparation parfaite des maisons.
\end{frame}

% ---------------------------------------------------------------------------
\coursepart{Calculer le gradient}{Gradient}{coursegradient}

\begin{frame}[t]{Le gradient est une moyenne de résidus pondérés}
  \vspace{4mm}
  \[
    \frac{\partial\ell_i}{\partial z_i}=p_i-y_i.
  \]
  Comme \(z_i=w_0+w_{\mathrm{pot}}x_{i,\mathrm{pot}}
                   +w_{\mathrm{vol}}x_{i,\mathrm{vol}}\),
  \[
  \boxed{
    \nabla J(w)=\frac1n\sum_{i=1}^{n}(p_i-y_i)
      \begin{pmatrix}1\\x_{i,\mathrm{pot}}\\x_{i,\mathrm{vol}}\end{pmatrix}
    =\frac1nX^\top(p-y).}
  \]
\end{frame}

\begin{frame}[t]{Gradient au point d'initialisation}
  \[
    w^{(0)}=(0,0,0)\quad\Longrightarrow\quad z_i=0,
    \quad p_i=\frac12.
  \]
  Les douze observations donnent exactement
  \[
    \sum_i(p_i-y_i)=0,\qquad
    \sum_i(p_i-y_i)x_{i,\mathrm{pot}}=3,\qquad
    \sum_i(p_i-y_i)x_{i,\mathrm{vol}}=-3.
  \]
  Donc
  \[
    \boxed{\nabla J(w^{(0)})=(0,\;0{,}25,\;-0{,}25).}
  \]
\end{frame}

\begin{frame}[t]{Contribution de quatre observations au gradient}
  \centering
  \renewcommand{\arraystretch}{1.22}
  \begin{tabular}{lrrrrr}
    \toprule
    élève & $y$ & $p-y$ à $w=0$ & $x_{\mathrm{pot}}$ &
      $(p-y)x_{\mathrm{pot}}$ & $(p-y)x_{\mathrm{vol}}$\\
    \midrule
    Harry   & $1$ & $-0{,}5$ & $-1{,}75$ & $+0{,}875$ & $+0{,}125$\\
    Drago   & $0$ & $+0{,}5$ & $-0{,}25$ & $-0{,}125$ & $-0{,}875$\\
    Neville & $1$ & $-0{,}5$ & $+1{,}00$ & $-0{,}500$ & $+0{,}250$\\
    Vincent & $0$ & $+0{,}5$ & $-0{,}50$ & $-0{,}250$ & $+0{,}500$\\
    \bottomrule
  \end{tabular}
  \vfill
  Les contributions sont additionnées avant la division par \(n=12\).
\end{frame}

\begin{frame}[t]{Calcul vectoriel du gradient}
  \vspace{3mm}
  \[
    \begin{aligned}
      z&=Xw,\\
      p&=\sigma(z),\\
      r&=p-y,\\
      \nabla J(w)&=\frac1{12}X^\top r.
    \end{aligned}
  \]
  \vfill
  \begin{coursecallout}{coursegradient!88!ink}{Dimensions}
    \(X\in\mathbb R^{12\times3}\), \(w\in\mathbb R^3\),
    \(p,y,r\in\mathbb R^{12}\) et
    \(\nabla J(w)\in\mathbb R^3\).
  \end{coursecallout}
\end{frame}

\begin{frame}[t]{Vérification du gradient par différence finie}
  \[
    \frac{\partial J}{\partial w_j}\approx
      \frac{J(w+h e_j)-J(w-h e_j)}{2h},\qquad h=10^{-5}.
  \]
  \centering
  \scriptsize
  \renewcommand{\arraystretch}{1.2}
  \begin{tabular}{crrr}
    \toprule
    composante $j$ & gradient analytique & différence finie & écart\\
    \midrule
    \MicroGradientCheckRows
    \bottomrule
  \end{tabular}
  \vfill
  Écart maximal : \(\MicroGradientMaxError\).
\end{frame}

% ---------------------------------------------------------------------------
\coursepart{Estimer les coefficients}{Mise à jour}{courseupdate}

\begin{frame}[t]{Première mise à jour des coefficients}
  \[
    w^{(t+1)}=w^{(t)}-\alpha\nabla J(w^{(t)}),\qquad \alpha=1.
  \]
  À partir de \(w^{(0)}=(0,0,0)\) :
  \[
    w^{(1)}=(0,0,0)-(0,0{,}25,-0{,}25)
             =\boxed{(0,-0{,}25,+0{,}25)}.
  \]
  Pour Harry, \(x_{\mathrm{vol}}-x_{\mathrm{pot}}=1{,}5\) :
  \[
    z^{(1)}=0{,}375,\qquad p^{(1)}=0{,}593.
  \]
  La perte moyenne passe de \(0{,}693147\) à \(0{,}585623\).
\end{frame}

\begin{frame}[t]{Convergence de la descente de gradient}
  \centering
  \scriptsize
  \renewcommand{\arraystretch}{1.08}
  \begin{tabular}{rrrrrr}
    \toprule
    itération & $J$ & $w_0$ & $w_{\mathrm{pot}}$ & $w_{\mathrm{vol}}$ &
      $\lVert\nabla J\rVert$\\
    \midrule
    \MicroGDTableRows
    \bottomrule
  \end{tabular}
  \vfill
  Le critère \(\lVert\nabla J\rVert\leq10^{-12}\) est atteint à l'itération
  \(\MicroGDIterations\).
\end{frame}

\begin{frame}[t]{Expression analytique de l'optimum}
  Les six observations vérifiant \(V>3P+20\) comprennent cinq Gryffondor et un
  Serpentard. Une probabilité commune \(q\) y maximise la vraisemblance pour
  \[
    q=\frac56.
  \]
  Leur différence standardisée vaut
  \(x_{\mathrm{vol}}-x_{\mathrm{pot}}=18/12=3/2\). En posant
  \(w_{\mathrm{pot}}=-a\) et \(w_{\mathrm{vol}}=a\),
  \[
    \sigma\!\left(\frac32a\right)=\frac56
    \Longrightarrow
    \frac32a=\operatorname{logit}\!\left(\frac56\right)=\ln5.
  \]
  \[
    \boxed{a=\frac{2\ln5}{3}=1{,}072959.}
  \]
\end{frame}

\begin{frame}[t]{Paramètres standardisés et paramètres bruts}
  \begin{columns}[T,onlytextwidth]
    \column{.48\textwidth}
      \begin{coursecallout}{courseupdate!88!ink}{Variables standardisées}
        \[
          (w_0,w_{\mathrm{pot}},w_{\mathrm{vol}})=\MicroW
        \]
        \[
          z=\frac{2\ln5}{3}
          (x_{\mathrm{vol}}-x_{\mathrm{pot}}).
        \]
      \end{coursecallout}
    \column{.48\textwidth}
      \begin{coursecallout}{coursescore!88!ink}{Notes brutes}
        \[
          \begin{gathered}
            b=-\frac{10\ln5}{9},\quad
            \beta_{\mathrm{pot}}=-\frac{\ln5}{6},\\
            \beta_{\mathrm{vol}}=+\frac{\ln5}{18}.
          \end{gathered}
        \]
        \[
          \boxed{z=\frac{\ln5}{18}(V-3P-20).}
        \]
      \end{coursecallout}
  \end{columns}
  \vfill
  \[
    J(w^\star)=\MicroTrainLoss.
  \]
\end{frame}

\begin{frame}[t]{Le recouvrement rend l'optimum fini}
  \centering
  \begin{tikzpicture}
    \begin{axis}[width=.74\textwidth,height=4.8cm,xmin=0,xmax=2.6,
      ymin=.43,ymax=.72,axis lines=box,grid=major,grid style={gridgray!30},
      xlabel={$a$ dans $(0,-a,+a)$},ylabel={$J(a)$},
      tick label style={font=\scriptsize},label style={font=\scriptsize}]
      \addplot[courseupdate!88!ink,very thick,domain=0:2.6,samples=120]
        {(5/6)*ln(1+exp(-1.5*x))+(1/6)*ln(1+exp(1.5*x))};
      \addplot[only marks,mark=*,mark size=2.8pt,warning]
        coordinates {(1.072958608,.450561209)};
      \node[font=\scriptsize,anchor=south west,text=warning]
        at (axis cs:1.072958608,.450561209) {$a^\star$};
    \end{axis}
  \end{tikzpicture}
  \begin{releve}
    Sans Neville et Vincent, les maisons seraient parfaitement séparées : la
    perte décroîtrait vers zéro lorsque \(a\) augmente, sans optimum fini.
  \end{releve}
\end{frame}

\begin{frame}[t]{Frontière du modèle ajusté}
  \centering
  \begin{tikzpicture}
    \begin{axis}[rawcloud,legend columns=2]
      \addplot[only marks,mark=square*,mark size=2.6pt,coursescore!88!ink]
        coordinates {\MicroTrainGryffondorCoords};
      \addlegendentry{Gryffondor}
      \addplot[only marks,mark=*,mark size=2.6pt,warning!88!ink]
        coordinates {\MicroTrainSerpentardCoords};
      \addlegendentry{Serpentard}
      \addplot[ink!82,very thick,domain=0:20,samples=2] {3*x+20};
      \node[font=\scriptsize\bfseries,text=warning,anchor=south east]
        at (axis cs:14,44) {Neville};
      \node[font=\scriptsize\bfseries,text=warning,anchor=north west]
        at (axis cs:8,62) {Vincent};
      \node[font=\scriptsize,fill=white,inner sep=1pt,text=ink!72]
        at (axis cs:15.5,64) {$p=0{,}50$};
    \end{axis}
  \end{tikzpicture}
\end{frame}

% ---------------------------------------------------------------------------
\coursepart{Définir la règle de décision}{Décision}{coursedecision}

\begin{frame}[t]{De la probabilité à la maison prédite}
  \vspace{4mm}
  \[
    \widehat y_i=\mathbf1[p_i\geq\tau].
  \]
  \begin{columns}[T,onlytextwidth]
    \column{.47\textwidth}
      \begin{coursecallout}{warning!88!ink}{$p_i<\tau$}
        Maison prédite : Serpentard.
      \end{coursecallout}
    \column{.47\textwidth}
      \begin{coursecallout}{coursescore!88!ink}{$p_i\geq\tau$}
        Maison prédite : Gryffondor.
      \end{coursecallout}
  \end{columns}
  \vfill
  Le seuil agit après l'estimation : le modifier ne recalcule ni la
  standardisation, ni les coefficients.
\end{frame}

\begin{frame}[t]{Probabilités produites sur le groupe d'estimation}
  \[
    V-3P-20=+18
    \Rightarrow z=+\ln5
    \Rightarrow p=\frac56,
  \]
  \[
    V-3P-20=-18
    \Rightarrow z=-\ln5
    \Rightarrow p=\frac16.
  \]
  \vspace{3mm}
  \begin{columns}[T,onlytextwidth]
    \column{.47\textwidth}
      \begin{coursecallout}{rulecolor!88!ink}{Dix observations correctement classées}
        \[
          \ell=-\ln(5/6)=0{,}182322.
        \]
      \end{coursecallout}
    \column{.47\textwidth}
      \begin{coursecallout}{warning!88!ink}{Neville et Vincent}
        \[
          \ell=-\ln(1/6)=1{,}791759.
        \]
      \end{coursecallout}
  \end{columns}
\end{frame}

\begin{frame}[t]{Résultats sur le groupe d'estimation}
  \centering
  \scriptsize
  \renewcommand{\arraystretch}{1.02}
  \begin{tabular}{cllrrrr}
    \toprule
    ID & élève & maison obs. & score & prob. & préd. 0,50 & préd. 0,65\\
    \midrule
    \MicroTrainCalcRows
    \bottomrule
  \end{tabular}
\end{frame}

\begin{frame}[t]{Matrice de confusion sur le groupe d'estimation}
  \centering
  \begin{tikzpicture}[x=1.45cm,y=.85cm,
    cell/.style={draw=white,line width=1pt,minimum width=2.8cm,
      minimum height=1.35cm,font=\Large\bfseries},
    lab/.style={font=\small\bfseries,text=ink!75}]
    \node[lab] at (1,3.25) {Serpentard prédit};
    \node[lab] at (3,3.25) {Gryffondor prédit};
    \node[lab,rotate=90] at (-.35,2.25) {Serpentard obs.};
    \node[lab,rotate=90] at (-.35,.75) {Gryffondor obs.};
    \node[cell,fill=rulecolor!18,text=rulecolor!82!ink] at (1,2.25) {VN $=5$};
    \node[cell,fill=warning!18,text=warning!88!ink] at (3,2.25) {FP $=1$};
    \node[cell,fill=warning!18,text=warning!88!ink] at (1,.75) {FN $=1$};
    \node[cell,fill=rulecolor!18,text=rulecolor!82!ink] at (3,.75) {VP $=5$};
  \end{tikzpicture}
  \vfill
  \[
    \text{exactitude}=\frac{5+5}{12}=83{,}3\%.
  \]
\end{frame}

\begin{frame}[t]{Le seuil déplace la frontière sans réestimer le modèle}
  \[
    p\geq\tau
    \Longleftrightarrow
    V\geq3P+20+\frac{18}{\ln5}\operatorname{logit}(\tau).
  \]
  \centering
  \begin{tikzpicture}
    \begin{axis}[width=.72\textwidth,height=4.75cm,
      xmin=0,xmax=20,ymin=0,ymax=100,
      xtick={0,5,10,15,20},ytick={0,20,40,60,80,100},
      grid=major,grid style={gridgray!32},axis lines=box,
      xlabel={Potions (sur 20)},ylabel={Vol (sur 100)},
      tick label style={font=\tiny},label style={font=\scriptsize},
      legend style={draw=none,font=\scriptsize,at={(.5,1.02)},anchor=south,legend columns=2}]
      \addplot[only marks,mark=square*,mark size=2.2pt,coursescore]
        coordinates {\MicroTrainGryffondorCoords};
      \addplot[only marks,mark=*,mark size=2.2pt,warning]
        coordinates {\MicroTrainSerpentardCoords};
      \addplot[ink!80,thick,domain=0:20,samples=2] {3*x+20};
      \addlegendentry{$\tau=0{,}50$}
      \addplot[coursedecision!88!ink,dashed,thick,domain=0:20,samples=2]
        {3*x+26.92};
      \addlegendentry{$\tau=0{,}65$}
    \end{axis}
  \end{tikzpicture}
\end{frame}

% ---------------------------------------------------------------------------
\coursepart{Évaluer le modèle hors échantillon}{Évaluation}{courseevaluation}

\begin{frame}[t]{Les paramètres sont figés avant l'évaluation}
  \vspace{5mm}
  \centering
  \begin{adjustbox}{max width=\textwidth,center}
  \begin{tikzpicture}[node distance=5mm,
    box/.style={draw=ink!25,rounded corners=2pt,minimum height=16mm,
      align=center,font=\scriptsize,inner xsep=3mm},
    arr/.style={-{Latex[length=2mm]},thick,ink!60}]
    \node[box,fill=courseevaluation!14] (row) {une observation\\à évaluer};
    \node[box,fill=coursedata!15,right=of row] (impute)
      {$\widetilde P=10$ ; $\widetilde V=50$\\si une note manque};
    \node[box,fill=courseprobability!15,right=of impute] (scale)
      {$\mu=(10,50)$\\$s=(4,12)$};
    \node[box,fill=courseupdate!15,right=of scale] (model)
      {$z=(\ln5/18)(V-3P-20)$\\coefficients figés};
    \node[box,fill=courseevaluation!20,right=of model] (out)
      {$p$, maison prédite, perte};
    \draw[arr] (row) -- (impute);
    \draw[arr] (impute) -- (scale);
    \draw[arr] (scale) -- (model);
    \draw[arr] (model) -- (out);
  \end{tikzpicture}
  \end{adjustbox}
  \vfill
  Aucune statistique et aucun coefficient ne sont réestimés avec les huit élèves.
\end{frame}

\begin{frame}[t]{Prédictions sur le groupe d'évaluation}
  \centering
  \scriptsize
  \renewcommand{\arraystretch}{1.11}
  \begin{tabular}{cllrrrr}
    \toprule
    ID & élève & maison obs. & score & prob. & préd. 0,50 & préd. 0,65\\
    \midrule
    \MicroTestCalcRows
    \bottomrule
  \end{tabular}
  \vfill
  Perte moyenne : \(J_{\mathrm{eval}}=\MicroTestLoss\).
\end{frame}

\begin{frame}[t]{Analyse locale autour du seuil de décision}
  \centering
  \begin{tikzpicture}
    \begin{axis}[rawcloud,legend columns=2]
      \addplot[only marks,mark=square,mark size=3pt,very thick,coursescore!88!ink]
        coordinates {\MicroTestGryffondorCoords};
      \addlegendentry{Gryffondor — évaluation}
      \addplot[only marks,mark=o,mark size=3pt,very thick,warning!88!ink]
        coordinates {\MicroTestSerpentardCoords};
      \addlegendentry{Serpentard — évaluation}
      \addplot[ink!78,very thick,domain=0:20,samples=2] {3*x+20};
      \node[font=\scriptsize\bfseries,text=warning,anchor=south east]
        at (axis cs:10,55) {Marcus : $p=\MicroProbabilityMarcus$};
      \node[font=\scriptsize\bfseries,text=coursescore!88!ink,anchor=north west]
        at (axis cs:13,65) {Lavande : $p=\MicroProbabilityLavande$};
    \end{axis}
  \end{tikzpicture}
\end{frame}

\begin{frame}[t]{Effet du seuil sur la matrice de confusion}
  \begin{columns}[T,onlytextwidth]
    \column{.48\textwidth}
      \centering
      {\bfseries\color{coursedecision!88!ink} $\tau=0{,}50$}\par\medskip
      \begin{tabular}{c|cc}
        & Serp. prédit & Gryff. prédit\\\hline
        Serp. obs. & $3$ & \cellcolor{warning!18}$1$\\
        Gryff. obs. & $0$ & \cellcolor{rulecolor!18}$4$
      \end{tabular}
      \par\medskip
      Marcus est classé à tort.
    \column{.48\textwidth}
      \centering
      {\bfseries\color{coursedecision!88!ink} $\tau=0{,}65$}\par\medskip
      \begin{tabular}{c|cc}
        & Serp. prédit & Gryff. prédit\\\hline
        Serp. obs. & \cellcolor{rulecolor!18}$4$ & $0$\\
        Gryff. obs. & \cellcolor{warning!18}$1$ & $3$
      \end{tabular}
      \par\medskip
      Lavande est alors classée à tort.
  \end{columns}
  \vfill
  L'exactitude reste \(7/8\) ; la nature de l'erreur change.
\end{frame}

\begin{frame}[t]{Comparaison des métriques selon le seuil}
  \centering
  \renewcommand{\arraystretch}{1.35}
  \begin{tabular}{rrrrrr}
    \toprule
    seuil & exactitude & précision & rappel & spécificité & F1\\
    \midrule
    \MicroTestSummaryRows
    \bottomrule
  \end{tabular}
  \vfill
  À \(\tau=0{,}50\) : précision \(=\MicroTestPrecision\),
  rappel \(=\MicroTestRecall\), spécificité \(=\MicroTestSpecificity\),
  \(F_1=\MicroTestFScore\).
\end{frame}

\begin{frame}[t]{Incertitude d'évaluation sur huit observations}
  \vspace{3mm}
  \[
    \widehat{\mathrm{acc}}=\frac78=87{,}5\%.
  \]
  L'intervalle de Wilson à 95\% vaut
  \[
    \boxed{[\MicroTestWilsonLow\,;\,\MicroTestWilsonHigh]}
    \approx[52{,}9\%\,;\,97{,}8\%].
  \]
  \begin{vigilance}[Portée de l'étude de cas]
    Les vingt observations sont synthétiques, équilibrées et choisies pour
    rendre chaque transformation vérifiable. Elles illustrent une méthode ;
    elles ne mesurent pas la performance attendue sur une population réelle.
  \end{vigilance}
\end{frame}



\gdef\coursechapterindicator{}
\gdef\coursepartshort{Ressources}
\gdef\coursepartposition{Ressources}
\gdef\coursepartindex{8}
\setbeamercolor{frametitle}{fg=white,bg=ink}
\begin{frame}[t,allowframebreaks]{Références}
  \printbibliography[heading=none]
\end{frame}

\end{document}
